\documentclass[letterpaper]{article}
\usepackage{iclr2027_conference,times}
\usepackage{amsmath,amssymb,graphicx,booktabs}
\usepackage{hyperref}
\usepackage{url}

% Official anonymous review mode: do not enable \iclrfinalcopy.
% Structural starter only. All content remains to be migrated and verified.
\title{RelaySpec: Adapting Frozen Speculative Drafters through Linear Feature Interfaces}
\author{Anonymous authors}
\newcommand{\draftnote}[1]{\par\noindent\textit{Drafting note: #1}\par}

\begin{document}
\maketitle
\begin{abstract}
Structural starter for section-by-section drafting. Replace this paragraph with the verified problem, approach, quantitative findings, and scope after completing the evidence audit. This document is not a submission-ready manuscript.
\end{abstract}

\section{Introduction}
\label{sec:introduction}
\draftnote{Problem, motivation, shared interface idea, two adaptation settings, and three supported contributions.}

\section{Related Work}
\label{sec:related}
\draftnote{Feature-conditioned drafting; draft adaptation and reuse; linear interfaces and structural re-parameterization.}

\section{Background and Problem Formulation}
\label{sec:problem}
\draftnote{Define target, reference model, drafter, feature interface, frozen parameters, and both adaptation settings.}

\section{RelaySpec}
\label{sec:method}
\subsection{Linear Feature Interface}
\draftnote{Feature maps, dimensions, normalization, initialization, and training/export overview figure.}
\subsection{Adaptation to Fine-Tuned Targets}
\draftnote{Target-generated supervision, attributed objective, and gradient flow.}
\subsection{Cross-Model Transfer}
\draftnote{Feature pairing, reconstruction objective, and tokenizer alignment.}
\subsection{Export and Inference}
\draftnote{Folding, runtime selection, cache semantics, and correctness assumptions.}

\section{Experiments}
\label{sec:experiments}
\subsection{Experimental Setup and Metrics}
\draftnote{Models, datasets, splits, selection, baselines, hardware, decoding, and metric denominators.}
\subsection{Adaptation to Fine-Tuned Targets}
\draftnote{Matched specialization results and strongest adaptation controls.}
\subsection{Cross-Model Transfer}
\draftnote{All transfer directions, workload variation, native recovery, and missing controls.}
\subsection{Shared Serving and Adaptation Cost}
\draftnote{Concurrency, per-request versus aggregate throughput, and measured cost with exclusions.}

\section{Analysis and Limitations}
\label{sec:analysis}
\subsection{Interface and Objective Controls}
\draftnote{Interpret matched evidence and identify remaining confounds.}
\subsection{Mechanistic Interventions}
\draftnote{Identity preservation, directional corrections, reconstruction interventions, and counterexamples.}
\subsection{Supervision and Cost Tradeoffs}
\draftnote{Data, anchors, coverage, checkpoint selection, and fitting cost.}
\subsection{Limitations}
\draftnote{Final verified scope, failure regimes, statistical limits, and reproduction gaps.}

\section{Conclusion}
\label{sec:conclusion}
\draftnote{Answer the central research question without introducing new claims.}

\subsection*{AI Use Statement}
\draftnote{Accurate disclosure of actual AI uses and author verification; also complete the submission form.}
\subsection*{Ethics Statement}
\draftnote{Relevant impacts, data provenance/privacy, and application limits.}
\subsection*{Reproducibility Statement}
\draftnote{Point to specific available recipes, code, manifests, and proof details.}

% Once verified citations have been added, replace the temporary heading with:
% \bibliography{references}
% \bibliographystyle{iclr2027_conference}
\section*{References}
\draftnote{Add verified author-year references using the official bibliography style.}

\clearpage
\appendix
\section{Experimental Reproducibility}
\section{Architecture, Alignment, and Verification}
\section{Complete Adaptation Controls}
\section{Supervision and Cost Results}
\section{Serving Results}
\section{Additional Mechanistic Evidence}
\section{Extended Related-Work Comparison}
\end{document}
